\documentclass[11pt]{article}
\usepackage[margin=1in]{geometry}
\usepackage{amsmath,amssymb}
\usepackage{mathtools}
\usepackage{booktabs}
\usepackage{microtype}
\usepackage[T1]{fontenc}
\usepackage{newtxtext,newtxmath}
\usepackage{xcolor}

\newcommand{\reac}{\mathrm{reac}}
\newcommand{\prodsup}{\mathrm{prod}}
\newcommand{\fs}{f_s^{\star}}
\newcommand{\code}[1]{\texttt{\small #1}}
\DeclarePairedDelimiter{\avg}{\langle}{\rangle}

\setlength{\parskip}{0.55em}
\setlength{\parindent}{0pt}

\title{\vspace{-1.5em}\textbf{\large \texttt{ray\_limit}: selectivity ray,
complete-reaction limit, and closure profile \(C(f)\)}}
\author{}
\date{}

\begin{document}
\maketitle
\vspace{-3.5em}
\begin{center}
\textcolor{gray!70!black}{\small Mixing-limited reaction closure \quad---\quad
\code{species\_limits.py}, \code{integrate\_species\_odes}}
\end{center}
\vspace{0.4em}

\section*{Setup}

Active species \(i\in\mathcal{A}\), reactions \(j=1,\dots,R\). Net stoichiometry matrix
\[
N\in\mathbb{R}^{|\mathcal{A}|\times R},\qquad
N_{ij}=\nu^{\prodsup}_{ij}-\nu^{\reac}_{ij}.
\]
Feed vectors \(Y^{(1)},Y^{(2)}\in\mathbb{R}^{|\mathcal{A}|}\) (stream~1 at \(f=1\),
stream~2 at \(f=0\)). The \textbf{no-reaction line} at mixture fraction \(f\in[0,1]\) is
\[
M_i(f)=f\,Y^{(1)}_i+(1-f)\,Y^{(2)}_i.
\]
Let \(y\) be the current ODE state, \(y^0\) its initial value, and \(\avg{\cdot}\) the
expectation under the step's \(\mathrm{Beta}(\alpha_t,\beta_t)\) pdf.

\section*{1.\quad Selectivity ray}

The accumulated change \(\Delta:=y-y^0\) lies in \(\operatorname{range}(N)\). The
\textbf{selectivity ray} is taken directly as this accumulated extent:
\[
d = \Delta = y - y^0.
\]
This preserves the true reaction direction without distortion from a nonnegativity clip
--- important when products form by more than one route (parallel reactions,
intermediates). Near \(y^0\) (where \(\Delta\approx 0\) and the direction is
ill-conditioned) the code falls back to the mass-action projection
\(d = N\,r(y)\), where \(r(y)\) is the vector of instantaneous mass-action rates.

The ray \(d\) encodes the \textbf{current selectivity}: which species have been consumed
or produced and in what ratio. It rotates as the ODE state evolves, tracking the shift
between competing pathways.

\section*{2.\quad The complete-reaction limit \(B(f)\)}

\(B_i(f)\) is built by reacting the mixing line \(M(f)\) to completion along the ray \(d\)
--- advancing as far as possible without any species going negative.

\textbf{Maximum extent.} At each \(f\) the largest non-negative scalar \(c_{\max}(f)\)
that keeps \(M(f)+c\,d\ge 0\) is
\[
c_{\max}(f) = \min_{\substack{i:\,d_i<0}}\frac{M_i(f)}{-d_i}.
\]
Only consumed species (\(d_i<0\)) constrain \(c_{\max}\); products (\(d_i>0\)) are
unconstrained. The limit is then
\[
B_i(f) = M_i(f) + c_{\max}(f)\,d_i.
\]

\section*{3.\quad Closure profile \(C_i(f)\) and scalar \(\lambda\)}

The closure profile interpolates between the complete-reaction limit \(B_i\) and the
no-reaction line \(M_i\):
\[
\boxed{\;C_i(f)=B_i(f)+\bigl[M_i(f)-B_i(f)\bigr]\,\lambda_i\;}
\]
\(\lambda_i=0\) gives the pure complete-reaction limit; \(\lambda_i=1\) gives pure
no-reaction mixing. \(\lambda_i\) is set so that \(\avg{C_i}=y_i\):
\[
\lambda_i=\operatorname{clip}_{[0,1]}\frac{y_i-\avg{B_i}}{\avg{M_i}-\avg{B_i}}.
\]
Segment integrals use the regularised incomplete Beta function \(I_x(\alpha,\beta)\)
(\code{betainc}):
\[
\int_{f_a}^{f_b}(sf+b)\,p_\beta\,df
= s\,\frac{\alpha_t}{\alpha_t+\beta_t}
  \bigl[I_{f_b}(\alpha_t{+}1,\beta_t)-I_{f_a}(\alpha_t{+}1,\beta_t)\bigr]
+ b\,\bigl[I_{f_b}(\alpha_t,\beta_t)-I_{f_a}(\alpha_t,\beta_t)\bigr].
\]

\textbf{Global \(\lambda\) for pure cross-stream.} Since \(M_i\) is linear,
\(\avg{M_i}=M_i(\bar{f})=y^0_i\), and \(B_i-M_i=c_{\max}\,d_i\) by construction, so
\[
\lambda_i = \frac{y_i-\avg{B_i}}{y^0_i-\avg{B_i}}
= \frac{d_i(1-\avg{c_{\max}})}{-\avg{c_{\max}}\,d_i}
= 1-\frac{1}{\avg{c_{\max}}}
\]
--- \textbf{the same scalar for every species} --- and \(\avg{C_i}=y_i\) exactly (no
clamp needed). Per-species \(\lambda_i\) diverge from this only when the clamp \([0,1]\)
binds.

\section*{4.\quad Worked illustration: \texttt{input\_fuller\_BC\_azo\_coupling\_kinetics}}

\textbf{Scheme:} R1: \(A{+}B\!\to\!R\), R2: \(A{+}R\!\to\!S\), R3: \(A{+}B\!\to\!T\),
R4: \(A{+}T\!\to\!S\), R5: \(A{+}C\!\to\!Q\).
\textbf{Feeds:} \(Y^{(1)}=15\,A\) (stream~1), \(Y^{(2)}=1.2\,B+1.2\,C\) (stream~2);
\(f_{\mathrm{mean}}=0.0625\).

\textbf{The effective single reaction.} The selectivity ray \(d=\Delta\) records the net
moles of each species consumed or produced per unit of accumulated reaction progress. Its
components \(d_i\) are the net stoichiometric coefficients of a single \emph{effective
lumped reaction} that encodes the current mixture of all five pathways:
\[
|d_A|\,A + |d_B|\,B + |d_C|\,C
\;\longrightarrow\;
d_R\,R + d_T\,T + d_S\,S + d_Q\,Q
\]
The coefficients rotate as the ODE state evolves. The table gives them at each
trajectory's output-grid time-midpoint (\(t=t_{\mathrm{end}}/2\)) for three dissipation
rates \(\varepsilon\) --- the actual \(A\) conversion reached at that instant is
\(77\%\), \(70\%\), and \(92\%\) respectively, not a fixed \(50\%\); the extent of
reaction still differs substantially at the same \emph{relative} point because the
trajectories are far from self-similar across \(\varepsilon\). At low \(\varepsilon\)
(broad pdf, persistent \(A\)-rich pockets) the secondary pathways dominate and \(S\)
carries most of the product weight; at high \(\varepsilon\) (near-uniform, \(A\)-lean)
the primary reactions dominate and \(R\) and \(Q\) accumulate.

\begin{center}\small
\begin{tabular}{l|ccc}
\toprule
 & \(\varepsilon=10^{-6}\) & \(\varepsilon=1\) & \(\varepsilon=10^{6}\)\\\midrule
\((\alpha_t,\beta_t)\) & \((0.311,\,4.66)\) & \((0.372,\,5.58)\) & \((14.5,\,218)\)\\
kink \(\fs\)           & \(0.1927\)         & \(0.1366\)         & \(0.0781\)\\
\(\lambda\) (global)   & \(0.000^{\dagger}\)& \(0.0008\)         & \(0.0541\)\\\midrule
\(d_A\) & \(-0.7225\) & \(-0.6556\) & \(-0.8616\)\\
\(d_B\) & \(-0.2422\) & \(-0.3314\) & \(-0.8135\)\\
\(d_C\) & \(-0.2421\) & \(-0.2868\) & \(-0.0470\)\\
\(d_R\) & \(+0.0040\) & \(+0.2832\) & \(+0.7560\)\\
\(d_T\) & \(+0.0000\) & \(+0.0108\) & \(+0.0564\)\\
\(d_S\) & \(+0.2382\) & \(+0.0374\) & \(+0.0011\)\\
\(d_Q\) & \(+0.2421\) & \(+0.2868\) & \(+0.0470\)\\
\bottomrule
\end{tabular}
\end{center}
{\small \(^{\dagger}\)At \(\varepsilon=10^{-6}\) the mean has reached the complete-reaction
average (\(y_i\ge\avg{B_i}\)), so \(\lambda\) clamps at \(0\) and
\(\avg{C_i}=\avg{B_i}\).}

\textbf{Building the limit (\(\varepsilon=1\)).} The effective reaction at this step is
\[
0.6556\,A + 0.3314\,B + 0.2868\,C \;\longrightarrow\; 0.2832\,R + 0.0108\,T + 0.0374\,S + 0.2868\,Q.
\]
The complete-reaction limit \(B(f)\) is found by advancing the mixing line \(M(f)\)
along this reaction as far as possible without a reactant going negative.  The maximum
extent \(c_{\max}(f)\) is set by whichever reactant is first exhausted:
\[
c_{\max}(f) = \min\!\left(\frac{M_A(f)}{0.6556},\;\frac{M_B(f)}{0.3314},\;\frac{M_C(f)}{0.2868}\right)
= \min\!\left(\frac{15f}{0.6556},\;\frac{1.2(1-f)}{0.3314},\;\frac{1.2(1-f)}{0.2868}\right).
\]
Since \(1.2(1{-}f)/0.3314 < 1.2(1{-}f)/0.2868\) for all \(f\), the binding constraints are
\(A\) and \(B\). Using \(p_i=M_i(0)/|d_i|\), \(q_i=(M_i(1)-M_i(0))/|d_i|\):
\[
(p,q)_A=(0,\,22.88),\quad (p,q)_B=(3.621,\,-3.621),\quad (p,q)_C=(4.183,\,-4.183).
\]
The \(A\)- and \(B\)-lines cross at the kink: \(\fs=3.621/26.50=0.1366\),
\(c^{\star}=3.127\). The peak values \(v^{\mathrm{fs}}_i=M_i(\fs)+c^{\star}d_i\) (e.g.\
\(v^{\mathrm{fs}}_R=0+3.127\times0.2832=0.886\)) complete \(B_i\); their Beta-averages
give \(\avg{B_i}\), and \(C_i=B_i+(M_i-B_i)\lambda\) with \(\lambda=0.0008\):

\begin{center}\small
\begin{tabular}{c|cccccc}
\toprule
\(\varepsilon=1\) & \(\avg{M_i}=y^0_i\) & \(v^{\mathrm{fs}}_i\) (mol\,m\(^{-3}\)) &
\(\avg{B_i}\) (mol\,m\(^{-3}\)) & \(y_i\) (mol\,m\(^{-3}\)) & \(\lambda_i\) &
\(\avg{C_i}\) \\\midrule
\(A\) & 0.9375 & 0.000 & 0.2814 & 0.2819 & 0.0008 & 0.2819 \\
\(B\) & 1.1250 & 0.000 & 0.7934 & 0.7936 & 0.0008 & 0.7936 \\
\(C\) & 1.1250 & 0.139 & 0.8379 & 0.8382 & 0.0008 & 0.8382 \\
\(R\) & 0.0000 & 0.886 & 0.2835 & 0.2832 & 0.0008 & 0.2832 \\
\(T\) & 0.0000 & 0.034 & 0.0108 & 0.0108 & 0.0008 & 0.0108 \\
\(S\) & 0.0000 & 0.117 & 0.0374 & 0.0374 & 0.0008 & 0.0374 \\
\(Q\) & 0.0000 & 0.897 & 0.2871 & 0.2868 & 0.0008 & 0.2868 \\
\bottomrule
\end{tabular}
\end{center}

Every \(\lambda_i\) equals the global scalar and \(\avg{C_i}=y_i\) to the quoted digits.

\section*{Note: single interior kink}

\textbf{Single interior kink.} For a pure cross-stream scheme the consumed species are
absent from their respective feed endpoints, so \(c_{\max}=0\) at \(f=0\) and \(f=1\).
Between the endpoints \(c_{\max}(f)\) is piecewise linear and concave; its kink \(\fs\)
is where the two tightest consumed-species constraints are simultaneously binding:
\[
\fs = \frac{p_b - p_a}{q_a - q_b},
\qquad p_i = \frac{M_i(0)}{-d_i},\quad q_i = \frac{M_i(1)-M_i(0)}{-d_i}.
\]
The peak extent at the kink is \(c^{\star}=c_{\max}(\fs)\). This gives the
\textbf{two-segment tent}:
\[
B_i(f) = \begin{cases}
M_i(f) + \dfrac{f}{\fs}\,c^{\star}\,d_i & f \le \fs\\[8pt]
M_i(f) + \dfrac{1-f}{1-\fs}\,c^{\star}\,d_i & f > \fs
\end{cases}
\]
with control-point peak value \(v^{\mathrm{fs}}_i = M_i(\fs)+c^{\star}d_i\) and
pure-feed endpoints \(B_i(0)=Y^{(2)}_i\), \(B_i(1)=Y^{(1)}_i\).

This structure holds when every species with \(d_i<0\) satisfies \(Y^{(1)}_i=0\) or
\(Y^{(2)}_i=0\).

\textbf{Intermediates are safe.} A species absent from both feeds has \(y^0_i=0\), and
since concentrations are non-negative, \(d_i=y_i-y^0_i=y_i\ge 0\). So an intermediate
can never have \(d_i<0\); the condition is vacuously satisfied for it.

\textbf{Condition failure does not cause collapse.} If a feed species appears in both
streams (\(Y^{(1)}_i>0\) and \(Y^{(2)}_i>0\)) and \(d_i<0\), then \(M_i(f)>0\) for
all \(f\) and \(c_{\max}\) remains positive. The only consequence is that \(B(f)\)
loses its endpoint-pinning (\(B(0)\ne Y^{(2)}\) or \(B(1)\ne Y^{(1)}\)); the
piecewise-linear machinery still computes the limit correctly.

\end{document}
